\documentclass{tietreport}

% ==================================================
% PACKAGES
% ==================================================

\usepackage{graphicx}
\usepackage{booktabs}
\usepackage{array}
\usepackage{amsmath}
\usepackage{url}
\usepackage{xcolor}
\usepackage{float}
\usepackage{enumitem}
\usepackage{listings}

% Bibliography configuration
\usepackage[
  date=year,
  style=alphabetic,
  backend=biber,
  sorting=nty,
  maxnames=5,
  minnames=3,
  maxalphanames=4,
  minalphanames=3,
  backref=true,
  doi=false,
  isbn=false,
  url=false,
  eprint=false
]{biblatex}

\addbibresource{references.bib}

% ==================================================
% CODE LISTING STYLE
% ==================================================

\lstset{
  basicstyle=\ttfamily\small,
  breaklines=true,
  frame=single,
  columns=fullflexible
}

% ==================================================
% TITLE PAGE
% ==================================================

\TitlePageHeader{%
  UCS503 - Software Engineering%
}

\ProjectTitle{%
  MemoryLane%
}

\ProjectSubTitle{%
  A Personal Web Memory and Semantic Search System%
}

\author{%
  \texttt{1024030716} Rabeeh Ahmad \\
  \texttt{1024030719} Ishan Singhania%
}

\TitlePageSubText{%
  Group: \texttt{3C51} BE Third Year, CSE%
}

\AdvisorName{%
  Dr. Raghav B. Venkataramaiyer%
}

\TitlePageFooterText{{%
    \large Mid-Semester Evaluation \\[0.2cm]
    \bfseries Computer Science and Engineering
    Department%
  } \\[0.15cm]
  Thapar Institute of Engineering and Technology,
  Patiala%
}

% ==================================================
% DOCUMENT
% ==================================================

\begin{document}

\maketitle

\tableofcontents

% ==================================================
% CHAPTER 1: INTRODUCTION
% ==================================================

\chapter{Introduction}

\section{Background and Context}

Modern web browsing involves interacting with a large amount
of information across articles, documentation, research
papers, tutorials, videos, and other online resources. A user
may encounter useful information at one point in time but
remember only a general idea of the content when they need it
again.

Conventional browser history provides a chronological list of
visited URLs, but it does not provide a content-aware way of
retrieving previously encountered information. Finding an old
resource therefore often requires remembering the website,
title, approximate date, or exact keywords that appeared on
the page.

MemoryLane is designed to address this problem by treating
previously encountered web content as a searchable personal
memory. Instead of simply recording that a webpage was visited,
the system stores its relevant content and creates a searchable
representation of that information.

A user can subsequently issue a natural-language query and
retrieve webpages whose stored content is relevant to the
query. The system supports both conventional lexical matching
and semantic similarity so that useful information can still be
found when the wording of the query differs from the wording
used on the original webpage.

\section{Problem Statement}

Traditional browser history is primarily concerned with
navigation rather than information retrieval. It records
information such as URLs and visit history, but it does not
provide an effective mechanism for answering content-based
questions about previously visited pages.

For example, a user may remember reading something about
``methods for reducing database query latency'' without
remembering the exact terminology or the website on which the
information appeared.

A personal web-memory system should therefore be able to:

\begin{enumerate}
    \item Capture and store useful information from visited
    webpages.

    \item Maintain webpage identity while processing long
    documents efficiently.

    \item Convert webpage content into a representation that
    supports meaningful search.

    \item Retrieve relevant information using natural-language
    queries.

    \item Handle both exact terminology and queries whose
    wording differs from the stored content.

    \item Present results as distinct webpages rather than
    exposing individual internal text chunks.

    \item Provide a practical interface between the browser,
    processing system, and search engine.
\end{enumerate}

\section{Project Objectives}

The primary objectives of MemoryLane are:

\begin{enumerate}

    \item \textbf{Personal information retrieval:}
    Develop a system that allows users to search previously
    encountered web content using natural-language queries.

    \item \textbf{Web content processing:}
    Build a pipeline for URL normalization, content validation,
    chunking, and representation of stored webpages.

    \item \textbf{Semantic representation:}
    Generate dense vector representations of webpage content
    to support meaning-based retrieval.

    \item \textbf{Efficient search:}
    Store processed webpage information in a search engine
    capable of both textual and vector-based retrieval.

    \item \textbf{Webpage-level results:}
    Ensure that multiple relevant sections of the same webpage
    are consolidated into a single user-facing result.

    \item \textbf{Evaluation:}
    Measure retrieval performance using established
    information-retrieval metrics and compare different search
    configurations.

\end{enumerate}

% ==================================================
% CHAPTER 2: METHODOLOGY AND DESIGN
% ==================================================

\chapter{Methodology and Design}

\section{System Architecture}

MemoryLane is organized as a pipeline connecting the browser,
user interface, backend processing layer, representation layer,
search index, and retrieval system.

The major components are:

\begin{enumerate}
    \item Browser extension
    \item Svelte frontend
    \item FastAPI backend
    \item URL normalization
    \item Content processing and chunking
    \item Sentence Transformer embedding model
    \item Apache Solr search index
    \item Retrieval and ranking module
    \item Optional summarization module
\end{enumerate}

The overall ingestion pipeline is:

\begin{center}
\textbf{
Browser Page / YouTube Video
$\rightarrow$
Browser Extension
$\rightarrow$
FastAPI
$\rightarrow$
URL Normalization
$\rightarrow$
Content Processing
$\rightarrow$
Chunking
$\rightarrow$
Embeddings
$\rightarrow$
Solr
}
\end{center}

The search pipeline is:

\begin{center}
\textbf{
User Query
$\rightarrow$
Svelte Frontend
$\rightarrow$
FastAPI
$\rightarrow$
Candidate Retrieval
$\rightarrow$
Ranking
$\rightarrow$
Webpage Results
}
\end{center}

The architecture separates data collection, presentation,
processing, storage, and retrieval. This makes individual
stages easier to develop, test, and modify independently.

\subsection{Browser Extension}

The browser extension serves as the collection layer of
MemoryLane.

When a user encounters a webpage, the extension captures
information associated with the page, including its URL, title,
and textual content. This information is sent to the backend
through the memory-ingestion API.

The extension does not perform the complete search or indexing
process itself. Instead, it acts as the interface between the
user's browsing environment and the MemoryLane backend.

\subsection{Svelte Frontend}

MemoryLane includes a lightweight frontend built using Svelte,
Vite, and JavaScript. The frontend provides the user-facing search
interface through which users can submit natural-language queries
and inspect retrieved memories.

The interface displays search results with information such as
page titles, URLs, and domains. It also supports on-demand
summarization of individual results.

The frontend communicates with the FastAPI backend through the
Vite development proxy, keeping the presentation layer separate
from the backend search and processing logic.

\subsection{FastAPI Backend}

The FastAPI backend provides the primary application interface
and coordinates communication between the browser extension,
frontend, and search infrastructure.

The main endpoints include:

\begin{itemize}
    \item \texttt{/health} --- checks the availability of
    required backend dependencies.

    \item \texttt{/memory} --- accepts webpage information and
    processes it for storage.

    \item \texttt{/search} --- accepts a natural-language query
    and returns relevant webpages.

    \item \texttt{/summarize} --- provides optional
    summarization functionality.
\end{itemize}

The backend is responsible for validating incoming data,
processing webpage content, generating embeddings, and
communicating with Solr.

\subsection{URL Normalization}

URLs can contain variations that still refer to the same
logical webpage. For example, differences in capitalization,
tracking parameters, URL fragments, default ports, or
\texttt{www} prefixes can result in multiple representations of
the same resource.

MemoryLane normalizes URLs before generating the webpage
identifier.

The normalization process includes:

\begin{itemize}
    \item Normalizing the URL scheme and hostname.
    \item Removing the \texttt{www} prefix.
    \item Removing default ports.
    \item Removing URL fragments.
    \item Removing common tracking parameters.
    \item Normalizing trailing slashes.
    \item Sorting query parameters.
\end{itemize}

A deterministic webpage identifier is generated from the
normalized URL. This identifier allows all chunks belonging to
one webpage to be associated with the same logical resource.

This is particularly important during retrieval because the
system ultimately presents webpages rather than individual
chunks.

\section{Content Processing}

\subsection{Content Validation}

Before content is indexed, the backend validates the submitted
URL and textual content. Extremely short or invalid content is
not treated as a useful webpage memory.

This prevents malformed or insufficient data from unnecessarily
occupying the search index.

\subsection{YouTube Content Handling}

MemoryLane also supports storing YouTube videos by using their
available transcripts as textual content. For a recognized YouTube
video URL, the backend attempts to retrieve the video's transcript
rather than treating the video player interface as webpage text.

The retrieved transcript is processed using the same pipeline as
normal webpage content. It is divided into chunks, converted into
dense representations, and stored in Solr. This allows information
contained within educational or informational videos to become
searchable alongside conventional webpages.

\subsection{Content Chunking}

Webpages can contain substantially more text than can be
effectively represented by a single vector. MemoryLane
therefore divides webpage content into smaller chunks.

The current configuration uses a minimum chunk size of
approximately 300 words and a maximum of approximately
500 words.

Chunking provides two important advantages. First, it allows
different sections of a long webpage to be represented
independently. Second, retrieval can identify the particular
section that is most relevant to a query.

Each chunk retains metadata connecting it back to its original
webpage.

The stored information includes:

\begin{itemize}
    \item Webpage identifier
    \item URL
    \item Domain
    \item Page title
    \item Chunk identifier
    \item Total number of chunks
    \item Timestamp
    \item Dense embedding
\end{itemize}

\subsection{Semantic Representation}

Each text chunk is converted into a dense vector using the
\texttt{all-MiniLM-L6-v2} Sentence Transformer model.

The model produces a 384-dimensional representation for each
chunk. The resulting vectors encode semantic characteristics
of the text and allow related concepts to be compared using
vector similarity.

Embeddings are normalized before similarity-based retrieval,
and cosine similarity is used for the dense-vector search.

\section{Information Retrieval}

MemoryLane uses more than one retrieval signal because
different queries benefit from different forms of matching.

\subsection{Lexical Retrieval}

The first retrieval mechanism is BM25-based lexical search.

BM25 is useful when a query contains specific words, names,
technical identifiers, or phrases that directly occur in the
stored content.

MemoryLane uses Solr's eDisMax query mechanism with the
\texttt{text} field as the primary field and the title field
given additional importance.

Lexical retrieval therefore provides a strong mechanism for
queries where the exact terminology is informative.

\subsection{Semantic Retrieval}

The second retrieval mechanism uses the dense representations
generated by the Sentence Transformer model.

When a user submits a query, the query is converted into the
same 384-dimensional embedding space as the stored chunks.
Solr then performs K-nearest-neighbour retrieval to identify
semantically similar content.

This allows MemoryLane to retrieve information when the user's
query and the original webpage use different words but describe
a similar concept.

\subsection{Combining Retrieval Signals}

The system can use the lexical and semantic retrieval signals
together.

The scores generated by the two retrieval methods are first
normalized independently using min-max normalization:

\[
s_i' =
\frac{s_i-s_{\min}}
{s_{\max}-s_{\min}}
\]

The current production configuration assigns equal weights to
the two signals:

\[
S =
0.5S_{\mathrm{BM25}}
+
0.5S_{\mathrm{vector}}
\]

This combined ranking is used by the main search pipeline,
while the individual retrieval configurations can also be
evaluated independently.

The purpose of combining the signals is not to replace either
retrieval method, but to provide the search system with
additional evidence when ranking results.

\subsection{Candidate Retrieval}

For a search query, the system retrieves a candidate pool from
the available retrieval mechanisms.

The current configuration retrieves up to 30 candidates from
each retrieval source before ranking the final results.

This candidate-generation stage allows the system to consider
a larger set of potentially relevant chunks than the number of
results ultimately presented to the user.

\subsection{Webpage-Level Result Aggregation}

The internal search process operates on chunks, but presenting
chunks directly would cause the same webpage to potentially
appear multiple times.

MemoryLane therefore groups retrieved chunks using their
\texttt{webpage\_id}.

For each webpage, the highest-scoring retrieved chunk is used
to represent that webpage's relevance. The webpages are then
ranked and returned as unique results.

The process can therefore be represented as:

\begin{center}
\textbf{
Chunk Retrieval
$\rightarrow$
Webpage Grouping
$\rightarrow$
Best-Chunk Selection
$\rightarrow$
Webpage Ranking
}
\end{center}

This maintains the advantages of fine-grained chunk retrieval
while producing a cleaner user-facing search experience.

\section{Technical Stack}

\begin{table}[H]
\centering
\begin{tabular}{p{0.32\textwidth} p{0.53\textwidth}}
\toprule
\textbf{Component} & \textbf{Technology} \\
\midrule
Programming language & Python \\
Frontend & Svelte / Vite / JavaScript \\
Backend framework & FastAPI \\
Search engine & Apache Solr 9.7 \\
Lexical retrieval & BM25 / Solr eDisMax \\
Vector retrieval & Solr KNN search \\
Embedding model & Sentence Transformers \texttt{all-MiniLM-L6-v2} \\
Embedding dimension & 384 \\
Containerization & Docker \\
Communication & HTTP / JSON \\
Optional summarization & Gemini-based summarization \\
\bottomrule
\end{tabular}
\caption{MemoryLane Technical Stack}
\label{tab:techstack}
\end{table}

% ==================================================
% CHAPTER 3: IMPLEMENTATION AND EVALUATION
% ==================================================

\chapter{Implementation and Evaluation}

\section{Implementation}

The backend is divided into modular components so that
different responsibilities remain separated.

The major implementation modules are:

\begin{itemize}

    \item \texttt{main.py} --- API endpoints and application
    orchestration.

    \item \texttt{models.py} --- request and response models.

    \item \texttt{url\_utils.py} --- URL normalization and
    webpage identifier generation.

    \item \texttt{retrieval.py} --- candidate retrieval,
    score processing, ranking, and webpage aggregation.

    \item \texttt{solr\_client.py} --- communication with the
    Apache Solr index.

    \item \texttt{config.py} --- system configuration and
    retrieval parameters.

\end{itemize}

\subsection{Memory Ingestion Pipeline}

When a webpage or supported YouTube video is submitted to the
\texttt{/memory} endpoint, the following sequence is performed:

\begin{enumerate}

    \item Validate the supplied URL.

    \item Normalize the URL.

    \item Generate a deterministic webpage identifier.

    \item Check whether the webpage already exists.

    \item Validate the supplied content.

    \item For supported YouTube URLs, retrieve the available
    transcript and use it as the textual content.

    \item Divide the content into chunks.

    \item Generate embeddings for the chunks.

    \item Remove outdated chunks associated with the same
    webpage when required.

    \item Store the processed chunks and metadata in Solr.

\end{enumerate}

This provides a consistent ingestion path for normal webpages,
supported YouTube videos, and evaluation data.

\subsection{Search Pipeline}

A search request follows the following sequence:

\begin{enumerate}

    \item Receive the natural-language query.

    \item Retrieve lexical candidates.

    \item Generate a dense representation of the query.

    \item Retrieve semantically similar candidates.

    \item Combine the available candidate results.

    \item Process their retrieval scores.

    \item Group chunks by webpage.

    \item Select the strongest representation of each webpage.

    \item Return the requested number of unique webpages.

\end{enumerate}

\section{Evaluation Methodology}

The retrieval component was evaluated using the SciFact
information-retrieval dataset.

The benchmark was adapted to the MemoryLane data model so that
the documents could pass through the actual ingestion pipeline.

The evaluation setup contained:

\begin{itemize}
    \item 1,000 unique documents
    \item 100 evaluation queries
    \item 106 relevant documents represented in the selected
    relevance data
\end{itemize}

The documents were indexed using the actual MemoryLane
\texttt{/memory} endpoint rather than being inserted directly
into Solr.

Therefore, the benchmark exercised the application's real
processing pipeline, including URL handling, content
processing, chunking, embedding generation, and indexing.

The benchmark documents use synthetic URLs of the form:

\begin{center}
\texttt{https://scifact.org/doc/\{document\_id\}}
\end{center}

These URLs provide unique webpage identifiers for the benchmark.
The actual document titles and scientific content originate
from SciFact.

Consequently, this experiment evaluates the information
retrieval component under a controlled corpus rather than
evaluating live web crawling or webpage extraction.

\section{Evaluation Metrics}

Three retrieval metrics were used to measure the quality of
the returned rankings.

\subsection{Mean Reciprocal Rank}

Mean Reciprocal Rank at ten (MRR@10) measures how early the
first relevant result appears among the top ten results.

For a query whose first relevant result occurs at rank $r$:

\[
RR = \frac{1}{r}
\]

The average reciprocal rank across all queries produces MRR.

A higher value indicates that relevant information tends to
appear earlier in the search results.

\subsection{Recall}

Recall@10 measures the proportion of relevant documents that
are retrieved within the first ten results:

\[
Recall@10 =
\frac{\text{Relevant documents retrieved in top 10}}
{\text{Total relevant documents}}
\]

A higher value indicates that fewer relevant documents are
missed by the retrieval system.

\subsection{Normalized Discounted Cumulative Gain}

nDCG@10 evaluates the quality of the ordering of the retrieved
results, giving greater importance to relevant results that
appear near the top.

Binary relevance was used for this experiment, meaning that a
document was treated as either relevant or non-relevant.

\section{Experimental Configurations}

Three retrieval configurations were evaluated:

\begin{enumerate}

    \item \textbf{BM25 Only:}
    Uses lexical retrieval without the semantic retrieval
    contribution.

    \item \textbf{Vector Only:}
    Uses dense semantic retrieval without the lexical
    contribution.

    \item \textbf{Combined Retrieval:}
    Uses the production configuration in which lexical and
    semantic retrieval signals are combined with equal
    weighting.

\end{enumerate}

Each configuration produced the top ten webpage-level results
for each of the 100 evaluation queries.

\section{Results}

The evaluation results are shown in Table~\ref{tab:evaluation}.

\begin{table}[H]
\centering
\begin{tabular}{lccc}
\toprule
\textbf{Method}
& \textbf{MRR@10}
& \textbf{Recall@10}
& \textbf{nDCG@10} \\
\midrule
BM25 Only
& 0.6482
& 0.8117
& 0.6776 \\

Vector Only
& \textbf{0.7795}
& 0.8867
& \textbf{0.8008} \\

Combined Retrieval
& 0.7538
& \textbf{0.9167}
& 0.7890 \\

\bottomrule
\end{tabular}
\caption{SciFact Retrieval Evaluation Results}
\label{tab:evaluation}
\end{table}

\section{Analysis of Results}

The results show that the different retrieval approaches have
different strengths on the selected benchmark.

\subsection{BM25 Retrieval}

BM25 achieved an MRR@10 of 0.6482, Recall@10 of 0.8117, and
nDCG@10 of 0.6776.

This provides a useful lexical baseline. Its performance
demonstrates that conventional term-based retrieval is capable
of retrieving a substantial portion of the relevant documents.

However, scientific queries may describe concepts using
language that differs from the exact wording in the relevant
documents. This can limit purely lexical retrieval.

\subsection{Semantic Retrieval}

Vector retrieval achieved the highest MRR@10, at 0.7795, and
the highest nDCG@10, at 0.8008.

This indicates that semantic representations were effective at
identifying relevant scientific content and placing relevant
documents near the top of the ranking.

The result is particularly relevant to MemoryLane because users
may remember the meaning of information they encountered without
remembering the exact terminology used on the original webpage.

\subsection{Combined Retrieval}

The combined configuration achieved the highest Recall@10,
with a value of 0.9167.

This indicates that incorporating both lexical and semantic
evidence allowed the system to retrieve a larger proportion of
the relevant documents within the top ten.

At the same time, vector-only retrieval achieved slightly higher
MRR@10 and nDCG@10. This shows that the best retrieval
configuration can depend on the characteristics of the
underlying documents and queries.

The current equal weighting therefore represents a practical
configuration rather than a claim that equal weighting is
universally optimal.

\section{Comparative Analysis}

The evaluation can be summarized as follows:

\begin{table}[H]
\centering
\begin{tabular}{p{0.30\textwidth} p{0.55\textwidth}}
\toprule
\textbf{Observation} & \textbf{Interpretation} \\
\midrule

Highest MRR@10 &
Vector retrieval places the first relevant result highest on
average. \\

Highest Recall@10 &
The combined configuration retrieves the largest proportion of
relevant documents within the top ten. \\

Highest nDCG@10 &
Vector retrieval provides the strongest ranking quality on the
selected benchmark. \\

BM25 baseline &
Lexical retrieval remains effective when query terminology
matches the stored content. \\

\bottomrule
\end{tabular}
\caption{Comparison of Retrieval Configurations}
\label{tab:comparison}
\end{table}

These results demonstrate why MemoryLane supports multiple
retrieval signals. Different search strategies provide
different forms of evidence, and their performance can vary
depending on the query and corpus.

% ==================================================
% CHAPTER 4: CONCLUSION
% ==================================================

\chapter{Conclusion}

\section{Summary of Work}

MemoryLane was developed as a personal web-memory system that
allows previously encountered web content to be stored,
processed, indexed, and retrieved through natural-language
search.

The system provides a complete pipeline beginning with browser
content collection and continuing through URL normalization,
content chunking, semantic representation, search indexing, and
result ranking.

The system also supports storing YouTube video content through
available transcripts, allowing video information to enter the
same processing and retrieval pipeline as conventional webpages.

The user-facing interface is implemented using Svelte, Vite,
and JavaScript, while the backend is implemented using FastAPI.
Apache Solr provides the underlying search infrastructure.
Webpage and transcript chunks are represented using
384-dimensional embeddings generated by the
\texttt{all-MiniLM-L6-v2} model.

The retrieval layer supports both conventional lexical search
and semantic vector search. These approaches can also be used
together to provide a broader set of retrieval evidence.

\section{Key Findings}

The evaluation produced several important findings.

\begin{enumerate}

    \item \textbf{MemoryLane successfully performs
    content-based retrieval.}
    The system was able to process and index a controlled corpus
    of 1,000 documents and evaluate retrieval over 100 queries.

    \item \textbf{Semantic representations were highly
    effective.}
    Vector retrieval achieved the highest MRR@10 of 0.7795 and
    nDCG@10 of 0.8008.

    \item \textbf{Multiple retrieval signals improved
    coverage.}
    The combined configuration achieved the highest Recall@10
    of 0.9167, showing that different retrieval mechanisms can
    complement one another.

    \item \textbf{Retrieval quality depends on the evaluation
    objective.}
    The configuration with the highest recall was not the same
    as the configuration with the highest MRR and nDCG. This
    highlights the importance of evaluating search systems using
    multiple metrics rather than relying on a single score.

    \item \textbf{Apache Solr provided an efficient search
    infrastructure.}
    The use of Apache Solr allowed lexical retrieval, vector
    retrieval, indexing, and ranking to be handled within a
    unified search platform, providing responsive retrieval
    during practical system testing.

\end{enumerate}

\section{Limitations}

The current implementation of MemoryLane has several project-level
limitations.

First, the quality of stored memories depends on the quality and
availability of the extracted webpage content. Dynamic webpages,
pages with limited accessible text, and YouTube videos without an
available transcript may not always provide sufficient content for
storage.

Second, the current system is primarily designed as a personal
web-memory prototype. Its storage and retrieval architecture has
not yet been extensively tested at the scale of very large
personal browsing histories.

Third, when the optional Gemini API is used for summarization,
personal browsing content must be sent to an external service.
This introduces privacy concerns because sensitive information
or identifiable browsing fingerprints could potentially be exposed.

Finally, the current user interface provides the core search and
result-viewing functionality, but additional usability features
would be required for a more complete history-management
experience.

\section{Future Work and Recommendations}

Several extensions can improve MemoryLane in future versions.

\begin{enumerate}

    \item \textbf{Improved retrieval ranking:}
    Investigate alternative ranking and score-combination
    strategies to improve search quality across different query
    types.

    \item \textbf{Improved chunking:}
    Explore sentence-aware, paragraph-aware, and overlapping
    chunking strategies.

    \item \textbf{Embedding models:}
    Compare the current embedding model with newer
    retrieval-oriented models.

    \item \textbf{Reranking:}
    Introduce a dedicated reranking stage for the strongest
    candidate results.

    \item \textbf{Larger-scale evaluation:}
    Test the system on larger and more diverse information
    retrieval datasets.

    \item \textbf{Scalability:}
    Study indexing and retrieval performance as the number of
    stored webpages increases substantially.

    \item \textbf{Recent browsing history integration:}
    Extend MemoryLane to maintain a recent history of accessed
    URLs alongside stored content memories. This would provide
    a more holistic browsing-history experience by combining
    conventional history features with content-aware retrieval.

    \item \textbf{Privacy and local processing:}
    Investigate local-first processing and storage approaches
    so that personal browsing information can remain under the
    user's control.

\end{enumerate}

\section{Final Remarks}

MemoryLane demonstrates how browser history can be extended
from a simple record of visited URLs into a searchable memory
of previously encountered information.

By storing webpage content, processing it into searchable
representations, and providing both lexical and semantic
retrieval capabilities, the system allows users to search for
information based on what they remember rather than only what
they can recall about the original webpage.

The evaluation demonstrates strong retrieval performance on
the selected scientific benchmark, while also showing that
different retrieval approaches have different strengths.
Overall, the project establishes a foundation for a practical
personal information-retrieval system that can be further
improved through better ranking, representation, scalability,
and real-world evaluation.

% ==================================================
% BIBLIOGRAPHY
% ==================================================

\printbibliography[heading=bibintoc]

\end{document}